\documentclass[10pt,a4paper,twocolumn]{article}
\usepackage[T1]{fontenc}
\usepackage{amsmath}
\usepackage{graphicx}
\usepackage{booktabs}
\usepackage{tabularx}
\usepackage{algorithm}
\usepackage{algpseudocode}
\usepackage{hyperref}
 
\title{Byte-Table Lookup Algorithm for Deploying Density-Based \\
       Classifiers on Memory-Constrained Devices}
\author{Andrei Prodan \\ Department of Computer Engineering Technology, Vanier College}
\date{25-09-2026}
 
\begin{document}
\maketitle
 
\section{Proposed Pipeline}

 Before beginning, this proposed pipeline does not involve the \emph{offline} process generated on a local system. Clustering (DBSCAN), kernel density estimation, normalization and building the byte tables are all assumed completed before ever reaching the \emph{online} process running on an external peripheral, which is our proposed pipeline that is solely inference and not training as well. The pipeline, which is logically demonstrated by Algorithm~\ref{alg:lookup}, requires certain inputs to function. The inputs are what was generated \emph{offline} mixed in with acquired data, meaning the pipeline uses a raw query point, the K byte tables, the scaling bounds, grid size G and threshold $\tau$. By using these inputs, a class label or UNKNOWN is outputted to classify the inputted point.
\begin{algorithm}[!bt]
\caption{Table-lookup classification (online path).}
\label{alg:lookup}
\begin{algorithmic}[1]
\Require A raw query point $(x,y)$
\Require One byte table $Q_c$ per class $c$, each $G \times G$, values $0$--$255$
\Require Scaling bounds $(x_{\min},x_{\max})$, $(y_{\min},y_{\max})$
\Require Grid size $G$ and reject threshold $\tau$
\Ensure A predicted class $\hat{c}$, or \textsc{unknown}
\State $\tilde{x} \gets \dfrac{x - x_{\min}}{x_{\max}-x_{\min}}$,\quad
       $\tilde{y} \gets \dfrac{y - y_{\min}}{y_{\max}-y_{\min}}$
\State $\text{col} \gets \lfloor \tilde{x}\,G \rfloor$, clipped to $[0,G-1]$
\State $\text{row} \gets \lfloor \tilde{y}\,G \rfloor$, clipped to $[0,G-1]$
\For{each class $c$}
    \State $v_c \gets Q_c[\text{row},\text{col}] / 255$ \Comment{one byte read}
\EndFor
\State $\hat{c} \gets \arg\max_c v_c$ \Comment{the winning class}
\If{$v_{\hat{c}} < \tau$}
    \State \Return \textsc{unknown} \Comment{outside every learned region}
\EndIf
\State $p_c \gets v_c / \sum_k v_k$ for each $c$ \Comment{optional: confidence}
\State \Return $\hat{c}$ (with distribution $p$)
\end{algorithmic}
\end{algorithm}
 
\subsection{Refinements}
Between Nearest-cell and Bilinear lookup, I would propose the use of Nearest-cell since it is 4x fewer flash reads per query. In a limited flash memory environment, consuming less memory is an inherent up-side, so between two options where confidence is more important than smoothness, choosing $K$ cost over $4K$ cost is a valid decision. Moreover, both options usually agree on the same outcomes, but may differ when dealing with edge cells.
As for the reject threshold, it is not optional because a value in a low enough part of the byte tables should not be allowed to be a specific class without being in an actually dense area. Also, for certain cluster shapes like the ring, no reject threshold could make it so points in the center could be considered a valid position to get the label, which goes against the fact the center of a ring cluster shouldn't be included in the cluster.

\section{Application}
 
A real-world application of this pipeline would be on a ATtiny85, which is a small microcontroller used on some Arduino boards. It has 8KB flash, 512 bytes SRAM and uses an 8-bit AVR core. Here are a few reasons why this chip would benefit from using the proposed pipeline:
\begin{itemize}
  \item \textbf{Flash budget}: The ATtiny85 has
        8192 bytes of flash total, of which the 768-byte model is roughly a 9.4 percent slice. Using any of the other byte table storage types would over-consume the limited storage, like float32 with a 37.5 percent slice or float64 with a 75 percent slice. Under those conditions, if a bootloader is required, it would be very hard to do anything with an extra 512 bytes being consumed.
  \item \textbf{Why no labelled training data}: It  takes too many resources to label each data point individually out in the field, so training using discovery would create organic connections that would mirror the real-world usage of the ATtiny85.
  \item \textbf{Why the classes wouldn't be round}: Let's say that you are collecting data with the ATtiny85, the temperature data has a tendency to curve from morning to night as the sun rises and sets. Most data gathering situations aren't round so why should the classes be as well? DBscan allows for classes to be the shape of the density cluster and not a center point like k-means which would cause these curved shapes to be treated as circles when classifying data.
\end{itemize}
 
\section{Method and Setup}
 
\begin{itemize}
  \item Dataset: The dataset being used is 1030 points, creating two crescents (340, 300) and
        a ring (300) plus 90 uniform noise points. It is synthetic data using the fixed
        seed \texttt{SEED = 20250903}
  \item Train/validation/test split: The dataset is split 60\%/20\%/20\%. The train set is what is used to generate the inference paths, with the validation set being use as reference during the training process to see how good it performs. The last set is the test set. The set is spent only once, the reason being that the test set can only be made up of data the system has never seen, else it would invalidate the test itself. So after running it once, the data cannot be reused to fulfill the same goal.
  \item Feature scaling: min--max to $[0,1]^2$ using training bounds only
  \item DBSCAN parameters: $\varepsilon = $ 0.06, $\text{min\_samples} = $
        8, with a cluster count of 3 (cluster 0: 209, cluster 1: 198, cluster 2: 178) and 33 noise points
  \item KDE bandwidth: $h = $ 0.05, selected on the
        \emph{validation} split by mean log-likelihood
        (Listing~5/\texttt{bandwidth\_sweep.py}) -- not eyeballed
  \item Grid size: $G = 16$; storage format: \texttt{uint8}
        (1 byte/cell); reject threshold: $\tau = 0.12$
\end{itemize}

\section{Evidence}

The evidence in this section follows the same order the pipeline was built in: the model is built using only the training split, the bandwidth $h$ is tuned using only the validation split, and the final result is measured once on the test split, which was not used for any earlier decision.
 
\subsection{Classes are discoverable without labels (Diagram B)}
Figure~\ref{fig:diagramB} shows DBSCAN and $k$-means on the same
training points. $k$-means uses a nearest-center rule, which unavoidably alters how the crescents and the ring are treated. Both of these elements would be treated like circles rather than their corresponding shapes, leading to point in the gaps where there is no data to be classified as included points instead of noise.
 
\begin{figure}[!bt]
\centering
\includegraphics[width=\columnwidth]{clustered_space.pdf}
\caption{Diagram B: DBSCAN preserves both crescents and the ring,
labelling background points as noise ($-1$); $k$-means with the
correct $k$ still slices the shapes apart. On the training split, DBSCAN recovered three clusters of 340, 300, and 300 points, correctly isolating the 90 scattered background points as noise, with no shape merged into another.}
\label{fig:diagramB}
\end{figure}
 
\subsection{The maps are well formed and survive storage as bytes (Diagrams C and D)}
Figure~\ref{fig:diagramC} and Figure~\ref{fig:diagramD} show the same probability surface, first as the smooth continuous version and then after it gets stored as bytes on the 16x16 grid. Storing it as uint8 gives a measured worst-case error of about 0.196\% of peak, which is small enough that you can't actually see it change the shape of any class.
 
\begin{figure}[!bt]
\centering
\includegraphics[width=\columnwidth]{total_distribution.pdf}
\caption{Diagram C: the continuous decision surface -- at every
location, the largest of the three normalised class maps, with the
reject plane at $\tau$. The surface dips below $\tau$ around the ring's hollow center and in the background area where there's no training data nearby, since those are the spots with no dense class close by.}
\label{fig:diagramC}
\end{figure}
 
\begin{figure*}[t]
\centering
\includegraphics[width=0.48\textwidth]{class_surfaces_continuous.pdf}
\hfill
\includegraphics[width=0.48\textwidth]{class_surfaces_discretized.pdf}
\caption{Diagram D: the per-class maps, continuous (left, Figure~3 of
the lab) vs.\ stored as \texttt{uint8} on the $16\times16$ grid
(right, Figure~5). The visible differences are smooth versus blocky. Fine map draws freely, while stored map is limited to 16x16. Quantization cannot be seen but the 0.195\% reduces below one shape shade.}
\label{fig:diagramD}
\end{figure*}
 
\subsection{What the model costs (Table~\ref{tab:footprint})}
Choosing uint8 is the best possible choice because it takes up the smallest amount of bytes while only having a worse-case error of 0.195\% from the actual peak.
 
\begin{table}[!bt]
\centering
\caption{Model footprint by storage format. $K=$ 3, $G=16$.}
\label{tab:footprint}
\begin{tabularx}{\linewidth}{@{}l r r X@{}}
\toprule
Storage & B/cell & Per class ($G^2$) & All $K$ classes \\
\midrule
\texttt{uint8}   & 1 & 256 & 768 \\
\texttt{uint16}  & 2 & 512 & 1536 \\
\texttt{float32} & 4 & 1024 & 3072 \\
\texttt{float64} & 8 & 2048 & 6144 \\
\bottomrule
\end{tabularx}
\\[2pt]
\footnotesize + 16 bytes of scaling bounds, shipped with every variant (2 floats per feature)
\end{table}
 
\subsection{One query resolved (Diagram E)}
Figure~\ref{fig:diagramE} shows a single held-out test point
classified via Algorithm~\ref{alg:lookup}.
The query point is [5.99958893,3.17868816]. It is located a cell [row11, col13]. The winning class/probability from my run was class 1 at a 100\% chance.
 
\begin{figure}[!bt]
\centering
\includegraphics[width=\columnwidth]{decision_map.pdf}
\caption{Diagram E: a held-out test query (star) in its grid cell,
and the per-class probabilities read from the stored tables. The
dashed line is $\tau$; the tallest bar is the arg max decision.}
\label{fig:diagramE}
\end{figure}
 
\subsection{Does the pipeline know when it doesn't know?}
When queried with the point [4.8, 1.6], which is located in the ring's center, the pipeline returned UNKNOWN. This answer matches the right answer for that point.
 
\subsection{Test-set result (headline number)}
On 206 held-out test points, the stored \texttt{uint8}
model agreed with the continuous KDE's own decision
100\% of the time. This means going from float64 down to uint8, an 8x reduction in memory, didn't cost any accuracy on data the model had never seen before, so uint8 is a safe pick for conserving memory.
 
\begin{table}[!bt]
\centering
\caption{Confusion between the continuous KDE decision (rows) and
the stored \texttt{uint8} decision (columns) on the test set.}
\label{tab:confusion}
\begin{tabularx}{\linewidth}{@{}l X X X X@{}}
\toprule
 & class 0 & class 1 & class 2 & unknown \\
\midrule
class 0 & 69 & 0 & 0 & 0 \\
class 1 & 0 & 68 & 0 & 0 \\
class 2 & 0 & 0 & 63 & 0 \\
unknown & 0 & 0 & 0 & 6 \\
\bottomrule
\end{tabularx}
\end{table}
 
\section{Limitations and Remaining Uncertainty}

This pipeline was only tested on one dataset, and that dataset is synthetic and only has two dimensions, so I can't say for sure it would behave the same way on real sensor data with more features.

The reject threshold $\tau = 0.12$ was a value I chose, not something the math derived on its own, so a different value could change how much of the grid gets rejected.

The 100\% test agreement is measured against what the continuous KDE itself decided, not against real ground-truth labels, because there aren't any -- DBSCAN is the one that invented the classes in the first place. So this number tells me the byte tables didn't lose anything compared to the KDE, but it can't tell me if the KDE's own decisions were "correct" in any deeper sense.

About half of the grid's cells fall below $\tau$ and get rejected, which means the model is being pretty conservative, only confidently classifying points that are actually close to where training data was seen.

If I doubled the grid to $G=32$, the footprint would go from 768 bytes to 3072 bytes ($32^2 \times 3$), a 4x jump in flash, while the resolution along each axis would only double. On a chip as small as the ATtiny85, that's likely not worth it unless the ring's hole specifically needed to be resolved better.

The bandwidth $h$ also matters a lot: too small and each class's map breaks up into separate little bumps around individual points; too large and the classes start merging into each other, which would shrink or fill in the ring's hole and risk it being classified as part of the ring instead of unknown.

Going forward, I'd want to test this pipeline on a real, labelled dataset instead of a synthetic one, so the test agreement number could be checked against actual ground truth instead of just the KDE's own decisions.
\end{document}